\chapter{記法と規約}
\label{app:conventions}
\chapterepigraph{``The last thing we decide on in writing a book is what shall be the first we put in it.''}{Attributed to Blaise Pascal, \textit{Pens\'ees}, Brunschvicg 19, trans.\ C.\ Kegan Paul}

本書で用いた規約をまとめる。異なる文献の公式を比較するときは、周波数の\term{時間規約}、
\term{Wronskian} の順序、\term{レゾルベント}の符号を最初に確認する必要がある。

\section{時空と単位系}

計量の符号は $(-,+,+,+)$、\term{幾何学単位系}は $G=c=1$ である。\term{静的球対称計量}を
\begin{equation}
\rmd s^2
=
-f(r)\,\rmd t^2
+f(r)^{-1}\rmd r^2
+r^2\rmd\Omega^2
\end{equation}
と書く。亀座標は $\rmd\rstar/\rmd r=f(r)^{-1}$ により定義する。

\section{\term{Fourier 規約}}

時間依存性は
\begin{equation}
\Psi(t,x)
=
\int
\frac{\rmd\omega}{2\pi}
\,
\rme^{-\rmi\omega t}
\psi(\omega,x)
\end{equation}
である。この規約では、$\im\omega<0$ が時間的減衰を表す。右向きの波は
$\rme^{-\rmi\omega(t-x)}$、左向きの波は $\rme^{-\rmi\omega(t+x)}$ である。

\section{演算子とレゾルベント}

\term{空間演算子}、\term{周波数演算子束}、\term{レゾルベント}を
\begin{align}
H
&=
-\frac{\rmd^2}{\rmd x^2}+V(x),
\\
\vL(\omega)
&=
H-\omega^2,
\\
\vR(z)
&=
(H-z)^{-1},
\\
\vG(\omega)
&=
\vL(\omega)^{-1}
\end{align}
と定めた。文献によっては $(z-H)^{-1}$ をレゾルベントと定義するため、\term{Riesz 射影}の
全体符号が変わる。本書の \term{contour moment} は式\eqref{eq:contour_moments}で直接定義し、
この曖昧さを避けた。

\section{Jost 解と Wronskian}

左右の \term{Jost 解}は
\begin{align}
f_-(x,\omega)
&\sim
\rme^{-\rmi\omega x},
&&x\longrightarrow-\infty,
\\
f_+(x,\omega)
&\sim
\rme^{+\rmi\omega x},
&&x\longrightarrow+\infty
\end{align}
で規格化する。\term{Wronskian} は
\begin{equation}
\vW(\omega)
=
f_-\frac{\partial f_+}{\partial x}
-
\frac{\partial f_-}{\partial x}f_+
\end{equation}
である。この順序に対して \term{Green 核}には式\eqref{eq:green_jost_formula}の負号が付く。

\section{略号}

本文では、\term{quasinormal mode} を QNM、\term{complex scaling method} を CSM、\term{exceptional
point} を EP と略記した。Regge--Wheeler と Zerilli は、それぞれ RW、Z と略記する
場合がある。

\section{量子二点関数}

遅延 Green 関数は
\begin{equation}
G_R(t)
=
\rmi\theta(t)\langle[O(t),O(0)]\rangle
\end{equation}
と定義する。Fourier 変換は $\int\rmd t\,\rme^{+\rmi\omega t}$ であり、この規約では
\term{スペクトル関数} $\rho_O=G^>-G^<$ と $\rho_O=2\im G_R$ が成り立つ。$\hbar=1$ とする。

\section{複素解析上の注意}

$z=\omega^2$ により、\term{エネルギー平面}と\term{周波数平面}は二重に覆われる。平方根、対数、
\term{Jost 解}の\term{解析接続}では、\term{切断}と \term{sheet} を固定する必要がある。数値計算では、複素平面
上の点だけでなく、そこへ至る\term{解析接続}の経路を記録することが望ましい。

例外点を周回するときは、制御パラメータの経路、固有ベクトルの\term{平行移動規約}、
\term{双直交内積}を指定する。\term{固有値の置換}、固有ベクトルの符号、\term{幾何学的位相}は関連するが、
同一の主張ではない。

\chapter*{改訂試作版について}
\addcontentsline{toc}{chapter}{改訂試作版について}

本改訂版は、「これから行う研究は研究問題へ送り、本筋に必要ないものは削る」という
方針で構成した。本文は、亀座標から古典的 Green 関数を作り、その極と連続成分を読み、
同じ遅延関数を量子交換子、KMS 条件、ETH の行列要素へ接続するところまでを一つの
主線とし、最後に応答不変量へ戻した。想定読者を学部四年生と明記し、未来地平面で
正則な座標、重力摂動の非放射部門、源から波形と flux への一つの完全な鎖、Kerr の
放射条件と \term{superradiance}、\term{面積エントロピー}と\term{第一法則}を物理ミニマムとして補った。
\term{Ryu--Takayanagi 公式}は本筋の導出には用いず、面積則から境界理論へ出るための短い
標識にとどめた。

通常の章末小問は置かない。極限地平面、長距離 Jost 関数、Kerr QNM の数値実装、
極とテールの一様近似、de~Sitter 連続応答、高次例外点、対称性を分解した ETH を
研究問題とした。
これらは本文の理解を確認する問題ではなく、本文の仮定を越えるための出発点である。

今後の改訂では、章を増やすこと自体を目標にしない。新しい結果が、源と観測者を結ぶ
因果的応答を変えるか、表示依存の量から新しい不変量を分離する場合にだけ本文へ戻す。
数値手法の比較、既知周波数の表、一般論の網羅は、主線を変えない限り研究問題または
別冊の実装ノートに置く。参考文献を除く本文百頁程度は編集上の目安であり、定義の順序を
壊す厳格な上限とはしない。
